\documentclass[10pt]{article}
\usepackage[margin=20mm]{geometry}
\usepackage{amsmath,amssymb,mathtools,fontspec,unicode-math,xurl,hyperref}
\IfFontExistsTF{DejaVu Serif}{\setmainfont{DejaVu Serif}}{\setmainfont{Latin Modern Roman}}
\setmathfont{latinmodern-math.otf}
\setlength{\parindent}{0pt}\setlength{\parskip}{6pt}\sloppy
\title{Technical Verification of the Generalised
André–Pink–Zannier Conjecture
for Shimura Varieties of Abelian Type}
\author{AI-assisted technical review}
\date{}
\begin{document}
\maketitle
Abstract: The acquired journal version is reviewed through its uniform integral Tate condition, Galois-orbit lower bound, integral norm comparison and independence arguments. A core proof-applicability gap is established: Proposition 5.7 invokes a comparison theorem requiring connected stabilizers, whereas its full centralizer stabilizer can remain disconnected at every good prime, including in a faithful adjoint representation. This is a localized Major finding about the derivation; it is not a counterexample to the principal theorem. The required fixed-centralizer comparison is repaired below by a finite-component quotient without changing its weight constant. Four auxiliary defects are separately assessed as Minor after tracing their actual uses. Exact counterexamples, counterevidence and repair boundaries are given below. Scientific adjudication remains pending.\par
\section*{Paper identity and inspected scope}
The reviewed paper is Rodolphe Richard and Andrei Yafaev, Generalised André-Pink-Zannier conjecture for Shimura varieties of Abelian type, Publications Mathématiques de l’IHÉS 141 (2025), 249–331, DOI 10.1007/s10240-025-00154-4 (https://doi.org/10.1007/s10240-025-00154-4) [paper]. The canonical input is the 83-page publisher version of record. A separately acquired 2023 arXiv source archive was not substituted for the journal text. Below, a PDF page is a one-based physical page; its printed page number is larger by 248.\par
The coverage ledger distinguishes full text reading, reconstructed arguments, actual selected image inspection and accepted external inputs. This is not a certification of every proof in the cited literature. No separate original supplement or computational repository was identified in the bounded availability checks. Lack of code is not a mathematical defect.\par
The principal strategy combines the uniform integral Tate property with a Galois-orbit lower bound and the earlier height/o-minimal argument. The exact published companion [height], rather than an assumed matching preprint, was consulted at the decomposition, reciprocity and main-proof interfaces. Its Lemma 6.12 corresponds to the older numbering cited in the reviewed paper.\par
\section*{Definitions and quantifiers that govern the review}
Definition 2.1 concerns a compact subgroup \(U\subset M(\widehat{\mathbb Z})\), with \(M\subset G\subset\operatorname{GL}_n\) connected reductive groups. The rational condition requires the appropriate connected Zariski closures, semisimplicity and centralizer equality. The residual condition is eventual and uniform in bounded index: for every fixed bound \(D\), a threshold \(M_U(D)\) controls all subgroups of the relevant residual image of index at most \(D\). It imposes no condition at every individual small prime.\par
The centralizer \(Z_G(M)\) is the full group scheme, not automatically its identity component. The abelianization \(M^{\mathrm{ab}}\), the full center \(Z(M)\), and the connected center are different objects. Likewise, stability under integral points \(G(\mathbb Z_p)\) is weaker than stability under the coaction of the integral group scheme \(G_{\mathbb Z_p}\). These distinctions are decisive in Problems 1, 2 and 4.\par
\section*{A disconnected stabilizer in the core norm-comparison invocation}
Problem 1 (Theorem 6.1 is invoked outside its connected-stabilizer hypotheses).  Proposition 5.7, PDF pp. 31–32 (printed pp. 279–280), uses Theorem 6.1 to prove equation (37). Its stabilizer is \(F=Z_G(M)\). Section 6.1, equation (49) and footnote 12, PDF p. 45 (printed p. 293), explicitly require the special fibers of both \(F\) and \(G\) to be connected reductive groups. Theorem 6.1 on the next page uses that setup and the full scheme stabilizers of its vectors. Definition 2.1 does not imply the required connectedness.\par
\subsection*{A model satisfying the operative group hypotheses}
Equation reference: eq:disconnected. Let \(W\) be a split four-dimensional quadratic space, put \(V=W\oplus\langle-1\rangle\), and take \[G=\operatorname{SO}(V)\simeq\operatorname{SO}_5,\qquad M=\operatorname{SO}(W)\simeq\operatorname{SO}_4\] with \(M\) fixing the last line. Over characteristic zero and every odd characteristic, the \(M\)-module \(W\) is absolutely irreducible and nontrivial. A commuting endomorphism of \(V\) therefore has the form \(\operatorname{diag}(aI_4,b)\). Preserving the quadratic form and determinant gives \[a^2=b^2=1,\qquad a^4b=1.\] Consequently \[
 Z_G(M)=\{\operatorname{diag}(I_4,1),\operatorname{diag}(-I_4,1)\}\simeq\mu_2.\] This has two geometric components at every odd prime.\par
Choose the product of the full integral local groups at good primes, with an arbitrary compact open integral subgroup at the finitely many remaining primes. Every local factor is Zariski dense in connected \(M_{\mathbb Q_p}\); the characteristic-zero representations are semisimple. For the residual clause, fix \(D\) and let \(q>\max(D,5)\). An index-at-most-\(D\) subgroup of \(M(\mathbb F_q)\) contains every element of order \(q\): the action on its cosets takes such an element to the identity in a symmetric group on fewer than \(q\) letters. It therefore contains \[\Omega_4^+(\mathbb F_q)
   =\operatorname{im}\bigl(\operatorname{SL}_2(\mathbb F_q)\times\operatorname{SL}_2(\mathbb F_q)
                          \longrightarrow \operatorname{SO}_4(\mathbb F_q)\bigr).\] The four-dimensional representation is the tensor product of the two standard two-dimensional representations; it is absolutely irreducible. Together with the trivial last line it is semisimple, and its commuting endomorphisms are exactly the two scalar blocks used above. Intersecting this commutant with \(G\) gives the same scheme-theoretic centralizer \(\mu_2\). This verifies the residual semisimplicity and centralizer conditions, uniformly for each fixed \(D\). Finally \(M^{\mathrm{ab}}=1\), so the additional abelian-image condition in Theorem 5.1 is automatic.\par
Passing to an adjoint representation does not remove the obstruction: \(G=\operatorname{SO}_5\) is adjoint and its adjoint representation is faithful. Moreover, \(\operatorname{Lie}G=\operatorname{Lie}M\oplus W\), and the nonidentity element of [eq:disconnected] (\#eq:disconnected) acts as \(-1\) on \(W\). The uniform Tate property passes through the faithful representation change used in the paper; alternatively, at large \(q\) the summands are the two simple three-dimensional adjoint modules and their four-dimensional tensor module, and the same residual check applies. Thus this obstruction to the literal connectedness hypothesis survives the adjoint image and exclusion of finitely many primes. The product group constructed here verifies the abstract theorem assumptions; it is not asserted to be the Galois image of a particular Shimura point.\par
\subsection*{Counterevidence and bounded impact}
The geometric auxiliary role also checked the related adjoint example \(\mathrm{PGL}_2\subset\mathrm{PGSp}_4\), obtained from two copies of the standard two-dimensional representation. Its centralizer is \(\mathrm{GO}_2/\mathbb G_m\), with two components. The detailed integral Tate verification is retained in the evidence files; the simpler model above already proves the missing implication.\par
\subsection*{A verified finite-component repair}
The fixed centralizer in this application admits the required extension after excluding finitely many primes. This keeps the full stabilizer \(F\), rather than substituting \(F^0\) in the orbit problem.\par
Put \(R=\mathbb Z_p\). Since \(F=Z_G(M)\) is fixed over \(\mathbb Q\), outside a finite set of primes its integral model is smooth, \(F^0_R\) is connected reductive, and \(Q=F_R/F^0_R\) is finite etale. Exclude also the prime divisors of its order \(m\). The connected quotient results in §7.3 and Lemma 7.11 [1, PDF pp. 56–58, 61] then give \[B=R[G]^{F^0_R},\qquad Y=\operatorname{Spec}B,
 \qquad Y_k=G_k/F^0_k\] for geometric fibers \(k\), with \(B\) of finite type and \(Y\) smooth. Right translation defines a \(Q\)-action on \(Y\). Set \[A=B^Q=R[G]^{F_R},\qquad X=\operatorname{Spec}A.\] After a finite etale base change making \(Q\) constant, the averaging operator \(m^{-1}\sum_{q\in Q}q\) is an idempotent onto the invariants. Thus invariants commute with passage to the residue and fraction fields. The orbit polynomial \(\prod_{q\in Q}(T-qb)\) shows that \(B\) is integral over \(A\); the finite-group invariant argument gives finite generation of \(A\) and finiteness of \(B\) over \(A\). These properties descend from the etale cover. Moreover \(A\subset R[G]\) is torsion-free, hence flat over the discrete valuation ring \(R\).\par
The action of \(Q_k\) on \(G_k/F^0_k\) is free: a component fixing \(gF^0_k\) must be the identity component. Its finite etale quotient is therefore smooth and \[X_k=(G_k/F^0_k)/Q_k=G_k/F_k .\] Finite presentation, flatness and geometrically smooth fibers imply that \(X/R\) is smooth. The same fiber criterion, applied to \(G\to X\), makes that morphism smooth and surjective. These are precisely the quotient inputs needed below.\par
For completeness, the normalization step also retains the full \(F\). Let \(S\) be the schematic orbit closure for a vector satisfying the other hypotheses of Theorem 7.1. In Lemma 7.10 [1, PDF pp. 59–61], \(X\) is normal, and \(X\to S\) is quasi-finite: its generic fiber is an isomorphism, and its special fibers have equal dimensions and the transitive orbit description just obtained. The normalization \(\widetilde S\) has a genuine \(G_R\)-action. Indeed the action lifts through normalization because \(G\times_R\widetilde S\) is normal by smooth base change; uniqueness gives the action law. Thus the corrected coaction hypothesis of Problem 4 applies. The proof of claim (3) of Lemma 7.10 uses Lemma 7.4, Corollary 7.5, stabilizer dimensions and connectedness of \(G\), with no further connectedness requirement on \(F\). It makes the special fiber of \(\widetilde S\) one geometric orbit. Zariski’s Main Theorem then identifies \(X\) with \(\widetilde S\), so \(X\to S\) is integral.\par
Proposition 7.15 [1, PDF p. 63] now gives the same integral-point criterion (67), over the integral closure \(\overline{\mathbb Z}_p\subset\overline{\mathbb Q}_p\) used in the paper. Smoothness and geometric surjectivity of \(G\to X\) also give the lifting step of Proposition 7.14. No assertion of surjectivity on \(\mathbb F_p\)-rational points is needed. In §6.3 [1, PDF p. 47], Cartan decomposition uses only the hyperspecial model of \(G\), and the two zero sets are exactly \[h_v(g)=0\quad\Longleftrightarrow\quad
 gF\in X(\overline{\mathbb Z}_p)
 \quad\Longleftrightarrow\quad h_{v'}(g)=0.\] The convex zero sets therefore coincide. Corollary 8.3 uses the unchanged representation weights, so the same constant \(C(\Sigma(\rho))\) works and no additive constant is introduced.\par
This finite exclusion is allowed: Proposition 5.7 states \(p\gg0\), and §5.1.3 [1, PDF p. 21] requires the uniform estimate only at almost all primes, with separate treatment of the remaining ones. The repair is verified for that actual fixed-centralizer application, using the stated connected quotient and normalization foundations. It is additional reasoning, not a claim that the authors supplied or adopted it.\par
\subsection*{Impact and assessment}
Equation (37) feeds Corollary 5.8 and the remaining local-height cases, then Theorems 5.4 and 5.1 and the Section 3 strategy for Theorem 1.4. The original invocation still fails an explicit hypothesis at this core step. Its assessment remains Major under the claim-level standard for a demonstrated invalid central derivation. The complete finite-component repair is recorded separately: the inequality and final theorem are not refuted, and the bridge is no longer an unresolved repair in this review.\par
\section*{The equality endpoint in the residual index assertion}
Problem 2 (Index \(p\) does not preserve the dagger subgroup).  Remark 2.1.3(2), PDF p. 5 (printed p. 253), asserts in paraphrase that \([U:U']\le p\) preserves the subgroup generated by residual unipotents: \(U'(p)^\dagger=U(p)^\dagger\). The equality endpoint is false.\par
Equation reference: eq:index. Fix a prime \(p_0\ge5\), take \(M=G=\operatorname{GL}_2\), and let \[C=\{I+aE_{12}:a\in\mathbb F_{p_0}\}.\] Set \(U_{p_0}\) equal to its full inverse image under reduction \(\operatorname{GL}_2(\mathbb Z_{p_0})\to\operatorname{GL}_2(\mathbb F_{p_0})\), and set \(U_q=\operatorname{GL}_2(\mathbb Z_q)\) at other primes. In the product \(U\), replace the \(p_0\)-factor by \(1+p_0M_2(\mathbb Z_{p_0})\) to obtain \(U'\). Then \[
 [U:U']=p_0,\qquad U(p_0)^\dagger=C,\qquad U'(p_0)^\dagger=1.\] All local factors are open and Zariski dense in \(\operatorname{GL}_2\). For a fixed \(D\), choose \(q>\max(p_0,D,3)\). The same coset-action argument forces every index-at-most-\(D\) residual subgroup to contain \(\operatorname{SL}_2(\mathbb F_q)\), which acts absolutely irreducibly with scalar commutant. Thus Definition 2.1 is satisfied; the exceptional prime \(p_0\) is permitted by its quantifiers.\par
The corrected hypothesis is \([U:U']<p\). A unipotent of \(p\)-power order acts trivially on fewer than \(p\) cosets, so all such elements lie in the smaller residual group and the dagger groups agree. This proves the corrected statement. Actual source images confirm the printed non-strict inequality; the example is not based on OCR.\par
The traced Section 5 uses have fixed index bounds followed by sufficiently large primes, so their cutoffs can impose \(p>D\); no essential use of the equality endpoint was found. The assessment is Minor, based on this dependency check and the proved strict-index assertion, not merely on the small size of the correction.\par
\section*{The subgroup constructed in Lemma 5.12}
Problem 3 (A kernel-intersection witness has a larger projective centralizer).  Lemma 5.12, PDF pp. 37–39 (printed pp. 285–287), seeks a bounded-index \(U'\subset U\) with \[Z_{M^{\mathrm{ad}}}(\operatorname{ad}U')=Z_M(U')/Z(M).\] Its proof specifies \(U[e]\), the intersection of the kernels of all homomorphisms \(U\to\mathbb Z/e\mathbb Z\), for a suitable \(e\). That specified subgroup need not satisfy the displayed equality.\par
Take \(M=\operatorname{SL}_2\) over \(\mathbb F_{17}\), \(e=2\), and \[g=\operatorname{diag}(2,9),\quad U=\langle g\rangle\simeq C_8,\quad
 h=g^2=\operatorname{diag}(4,13).\] Here \(e\) is divisible by the order of \(Z(M)(\overline{\mathbb F}_{17})\), and the prescribed subgroup is \(U[2]=\langle h\rangle\simeq C_4\). For \[w=\begin{pmatrix}0&-1\\1&0\end{pmatrix}\] direct calculation gives \(whw^{-1}=-h\). Thus \(\operatorname{ad}(w)\) centralizes \(\operatorname{ad}(U[2])\), although \(w\), and every centrally modified lift of it, fails to centralize \(h\). Since \(h\) has distinct eigenvalues, \(Z_M(U[2])\) is the diagonal torus; the projective centralizer has an additional Weyl component.\par
An independently written finite-field enumeration retained in `native\_galois\_cyclic\_calculation.json` checks the same matrix equalities and counts 16 elements in the \(\operatorname{SL}_2(\mathbb F_{17})\) centralizer of \(h\), against 32 in the inverse image of the projective centralizer. The direct matrix witness proves the discrepancy without identifying rational points of a quotient with a quotient of rational points.\par
Equation reference: eq:weaker. The commutator argument does prove the weaker statement \[
 Z_{M^{\mathrm{ad}}}(\operatorname{ad}U)
       \subseteq \operatorname{im}\bigl(Z_M(U[e])\to M^{\mathrm{ad}}\bigr).\] Indeed, for a lift of an element on the left, its commutators with \(U\) form a homomorphism into the finite central group; the kernel intersection kills its cyclic coordinates. This argument keeps \(U\) on the left. It does not permit replacing it there by \(U[e]\).\par
The existential assertion of Lemma 5.12 is not refuted by this example: \(U'=1\), of index 8, works in this individual case. Its general strong form remains unverified in this review. In the actual use in Lemma 5.11, however, [eq:weaker] (\#eq:weaker) is enough: the Tate property identifies the centralizer of the selected bounded-index subgroup with \(Z(M)\), forcing the projective centralizer of the original subgroup to be trivial. Thus the assessment is Minor: the proposed witness and stronger proof assertion fail, whereas the required downstream implication has a complete weaker replacement.\par
\section*{Integral-point stability versus a group-scheme action}
Problem 4 (Proposition 7.2 assumes too little stability).  Proposition 7.2, PDF pp. 51–52 (printed pp. 299–300), assumes a finitely generated \(G(\mathbb Z_p)\)-stable subalgebra of \(\mathbb Z_p[G]\) and concludes an equivariant integral representation/orbit realization. Pointwise stability does not supply the required scheme action.\par
Equation reference: eq:coaction. Let \(G=\mathbb G_m\) and \[A=\mathbb Z_p[pt,t^p-t]\subset\mathbb Z_p[t,t^{-1}].\] This meets the stated finite-generation and smooth reductive group hypotheses. For every \(u\in\mathbb Z_p^\times\), \[u(pt)=u\,pt,\qquad
 u(t^p-t)=u^p(t^p-t)+\frac{u^p-u}{p}\,pt ,\] so \(A\) is stable under all integral points. On the other hand, \(t\notin A\): modulo \(p\), the image of \(A\) is \(\mathbb F_p[t^p-t]\), whose nonconstant elements have degree divisible by \(p\). The scaling coaction satisfies \[
 \delta(t^p-t)=(t^p-t)s^p+t(s^p-s).\] Its coefficient of \(s\) is \(-t\notin A\), so \(\delta(t^p-t)\notin A[s,s^{-1}]\). The specified generic action cannot extend to this integral algebra. Normality or an equality of generic fibers is not a hypothesis of the printed proposition.\par
A sufficient repair is to require stability under the \(G_{\mathbb Z_p}\)-coaction. The actual applications have this stronger property. In Lemma 7.6, \(A=\mathbb Z_p[G]^F\); left and right translations commute, and flat base change preserves the kernel defining invariants, giving the genuine left coaction. In Lemma 7.10, the action on a schematic orbit closure lifts uniquely to its normalization: the product with the smooth group is normal, and the dominant action factors through the normalization. Uniqueness gives the group law. The normalization interfaces are supported by the precise primary statements in [stacks].\par
The assessment is Minor: the general auxiliary proposition is false, while the two traced applications satisfy the corrected hypothesis. This coaction repair is also used in the normalization step of the finite-component extension proved in Problem 1.\par
\section*{The conjugacy-finiteness prelude to Lemma B.5}
Problem 5 (Semisimple embeddings have infinitely many Frobenius twists).  The first preliminary bullet before Lemma B.5, PDF p. 72 (printed p. 320), asserts boundedly many \(\operatorname{GL}_n(\overline{\mathbb F}_p)\)-conjugacy classes of semisimple algebraic subgroups for fixed \(n\), uniformly in \(p\). Even at one fixed prime and dimension, the unrestricted assertion is false.\par
Fix \(p\ge5\), \(n=4\), and for every \(r\ge1\) let \[S_r=\{\operatorname{diag}(g,\operatorname{Frob}^{r}(g)):g\in\operatorname{SL}_2\}.\] This is an \(\mathbb F_p\)-defined closed connected semisimple subgroup, isomorphic to \(\operatorname{SL}_2\) by its first projection. Its Lie algebra is \(\operatorname{diag}(\mathfrak{sl}_2,0)\). The common fixed space of that Lie algebra is the second two-dimensional block, and the span of its images is the first block. Any conjugacy \(S_r\to S_s\) must therefore be block diagonal, say \(\operatorname{diag}(a,b)\). Its induced equations, applied to \(g=I+tE_{12}\), imply \[I+t^{p^r}N_1=I+t^{p^s}N_2\] for nonzero constant matrices \(N_1,N_2\). Polynomial degrees force \(r=s\). Thus the groups are pairwise nonconjugate even over the algebraic closure.\par
The actual printed conjugating field is \(\overline{\mathbb F}_p\); the review does not use the incorrect OCR reading \(\mathbb F_p\). There is also relevant mathematical counterevidence: \(S_r(\mathbb F_p)\) is always diagonal \(\operatorname{SL}_2(\mathbb F_p)\), so its Nori envelope is diagonal \(\operatorname{SL}_2\), not \(S_r\). The displayed family does not contradict the narrower Nori-group setting actually needed.\par
More directly, the two structural bounds used in Lemma B.5 do not need finiteness of embeddings. A semisimple subgroup of \(\operatorname{GL}_n\) has rank at most \(n\). There are finitely many semisimple root data of that bounded rank, including the finitely many central-isogeny types. Their centers have uniformly bounded order, and their outer automorphism groups are bounded by automorphisms of the corresponding based root data. Taking maxima gives the required constants. This proves the needed replacement without the false embedding-finiteness assertion.\par
The assessment is Minor, limited to an unnecessarily strong auxiliary assertion. No counterexample to Lemma B.5 or Theorem B.1 is established.\par
\section*{Other checks, overall impact and limits}
Section 4’s specialization and finite-type-field independence arguments were checked at the stated interfaces. The general-field independence input is supplied by the exact Gajda–Petersen result [gp]; a number-field-only objection is not retained. The auxiliary role consulted the relevant Noot and Serre passages [noot, serre], without claiming a full audit of all external proofs.\par
Several candidates were discharged. The torus in Section 6.1 can be chosen to fit the hyperspecial point; footnote 13 and the freedom of choice support that interpretation. An arbitrary badly positioned torus need not have a reductive closure, but that observation is not an additional finding against the required choice. In Lemma B.10, the \(e\)-power map has derivative \(e\,\mathrm{id}\) over \(\mathbb Q_p\). Its image contains an identity neighborhood, so a subgroup containing all \(e\)-th powers is open; taking its normal core makes the stated restricted-Burnside step applicable. The global projection/pure-factor argument can use the subsequent independent cyclic extraction argument in an earlier order. Actual images retain the closure bars lost by text extraction.\par
The isolated description of the full center as a torus in Appendix B.1 should read a smooth group of multiplicative type, with torus identity component: \(\operatorname{SL}_2\) has full center \(\mu_2\) at odd primes. Later steps already allow finite central parts. No additional substantive failure is assigned to that wording. Uniquely recoverable typesetting slips in nearby formulas are likewise not numbered Problems.\par
The supported overall assessment is Major, entirely on the core dependency in Problem 1. Problems 2–5 are local Minor findings for the explicit reasons given above. The review does not prove the main conclusion false. The finite-component repair of the disconnected-stabilizer interface is verified for the actual eventual-prime application, while the original invalid core invocation and its severity remain distinct from that repair. Masser–Wüstholz/Faltings, complete reducibility, Nori, Jordan, restricted Burnside and the o-minimal foundations remain accepted external inputs in their recorded scopes. Full access to the original Seshadri input was blocked in the bounded retrieval attempt; no broader theorem was inferred from that unavailable text. No matching author correction was identified in the bounded search; priority and author adoption are unknown.\par
\section*{Conclusion and separate decision states}
Item | Assessment | Repair verified / remaining work\par
Problem 1 | Major, invalid core invocation | Finite-component bridge proved for the actual eventual-prime application.\par
Problem 2 | Minor, false endpoint | Strict-index version proved; actual uses traced.\par
Problem 3 | Minor, invalid specified witness | Weaker required inclusion proved; stronger lemma unresolved.\par
Problem 4 | Minor, false generality | Coaction hypothesis restores the two actual applications.\par
Problem 5 | Minor, false embedding finiteness | Needed structural bounds proved from bounded root data.\par
Scientific adjudication is pending. Author adoption is unknown. No human review has been supplied. The present revision incorporates focused independent acceptance of the five existing Problems.\par
9 R. Richard and A. Yafaev, Generalised André-Pink-Zannier conjecture for Shimura varieties of Abelian type, Publ. Math. IHÉS 141 (2025), 249–331. doi:10.1007/s10240-025-00154-4 (https://doi.org/10.1007/s10240-025-00154-4). Reviewed publisher version. R. Richard and A. Yafaev, Height functions on Hecke orbits and the generalised André–Pink–Zannier conjecture, Compos. Math. 160 (2024), 2531–2584. doi:10.1112/S0010437X2400722X (https://doi.org/10.1112/S0010437X2400722X). Published interface passages consulted. The Stacks Project, Tag 035Q, Lemma 29.55.5(4) (https://stacks.math.columbia.edu/tag/035Q), and Tag 07TD, Lemma 37.19.3 (https://stacks.math.columbia.edu/tag/07TD). Normalization factorization and smooth base change; accessed 26 October 2026, not asserted to be a 2025 snapshot. W. Gajda and S. Petersen, Independence of \(\ell\)-adic Galois representations over function fields, Compos. Math. 149 (2013), 1091–1107, doi:10.1112/S0010437X12000711 (https://doi.org/10.1112/S0010437X12000711). Publisher abstract and the interface in arXiv:1103.2893v2 (https://arxiv.org/abs/1103.2893v2), Theorem 1.1 and Corollary 1.2, consulted by the Tate auxiliary role; journal equivalence of that author PDF is not asserted. R. Noot, Abelian varieties—Galois representation and properties of ordinary reduction, Compos. Math. 97 (1995), 161–171. Numdam PDF (https://www.numdam.org/item/CM\_1995\_\_97\_1-2\_161\_0.pdf); specialization interfaces in Section 1 consulted by the Tate auxiliary role. J.-P. Serre, Un critère d’indépendance pour une famille de représentations \(\ell\)-adiques, Comment. Math. Helv. 88 (2013), 1–12, item [47] in the reviewed paper’s bibliography. Author-hosted PDF (https://www.college-de-france.fr/media/jean-pierre-serre/UPL8570957518086251784\_independance.pdf); number-field setup, Sections 3.1–3.2 and the finite-type-field footnote consulted by the Tate auxiliary role.\par
\section*{Supporting checks}
The canonical PDF SHA256 is \[\texttt{3f3af58def7bb398b3d365ecf828a6571a9}\] \[\texttt{477fd7f2b7b891706f146fec05b9e}.\] Exact source/manifest bindings, the distinct arXiv-source identity, bounded proof checkpoints, calculation input/output and selected source-page images are retained in the per-paper verification directory. The finite calculation is reviewer-written; no author code was executed.\par
In the initial review, the root role integrated the science and wrote the report. They inherited the active runtime without a model override; their checks are not independent acceptance. No human identity or institutional endorsement is supplied. Scientific decision is pending and the human-review list is empty.\par
During independent acceptance, the primary session checked the five existing Problems and verified the finite-component repair with bounded read-only support from the native role `/root/artifact\_binding\_check`. That role recommended lowering Problem 1 to Minor after repair; the primary session retained Major because the demonstrated original invocation lies in the core derivation, while changing the repair state to verified. This revision preserves that distinction. The primary session performs the new report build and all final-page inspections; their actual execution and viewing facts are recorded separately in the verification files.\par
\end{document}
